\chapter*{Anexo A}
\section*{Encuesta de accesibilidad}
A continuación se listan las preguntas presentes en la encuesta de accesibilidad otorgada a los usuarios que realizaron las pruebas del binario final. Se presenta también el formato disponible para la respuesta y se cumplimenta un enlace \footnote{ \url{https://docs.google.com/forms/d/e/1FAIpQLSe8AECLV-pSyYLHRe6Y3eLFm-KL4B4qaLUY8uYthd5Ykz3l3Q/viewform?usp=sharing&ouid=115849952895919619315}} donde se publicó para su uso.
\begin{enumerate}
	\item[$\blacksquare$] Las siguientes preguntas están relacionadas con la jugabilidad y las interfaces del juego.
	\item ¿En qué medida has podido ubicarte en el mapa para localizar los puzles?
	\subitem- Respuesta del 1 al 10. 1 significa \say{Me he perdido y nunca me orienté}, 10 significa \say{Me he podido orientar perfectamente}.
	\item ¿Han sido las herramientas de ubicación (sonar y colisiones) suficientes y necesarias? Indique sus opiniones sobre ellas.
	\subitem- Texto de respuesta larga.
	\item ¿Los controles de teclado son intuitivos?
	\subitem- Sí, No, u otra respuesta de redacción libre.
	\item ¿Los controles de mando son intuitivos?
	\subitem- Sí, No, u otra respuesta de redacción libre.
	\item ¿Son los controles del juego compatibles con las tecnologías de asistencia instaladas en su ordenador?
	\subitem- Sí, No, u otra respuesta de redacción libre.
	\item ¿La forma en la que los menús están organizados es intuitiva y fácil de navegar?
	\subitem- Sí, No, u otra respuesta de redacción libre.
	\item ¿La narración por ordenador ha sido intuitiva o has echado de menos un narrador humano?
	\subitem- Texto de respuesta larga.
	\item[$\blacksquare$] Las siguientes preguntas están relacionadas con los diferentes modos de juego.
	\item ¿En qué medida valoras poder jugar con otras personas cooperativamente?
	\subitem- Texto de respuesta corta.
	\item ¿Consideras que un juego asimétrico (controles distintos para cada jugador) es un buen diseño para un juego cooperativo accesible? ¿En qué medida?
	\subitem- Texto de respuesta larga.
	\item ¿Has preferido la versión multijugador o la de un solo jugador? ¿Por qué?
	\subitem- Texto de respuesta corta.
	\item[$\blacksquare$] Las siguientes preguntas están relacionadas con las distintas mecánicas de sonificación.
	\begin{itemize}
		\item \textit{Preguntas sobre el tutorial.}
	\end{itemize}
	\item ¿Te ha parecido que esta mecánica de juego sonificada tiene potencial para ser utilizada en futuros juegos?
	\subitem- Respuesta del 1 al 10. 1 significa \say{Ha sido frustrante}, 10 significa \say{Ha sido divertida}.
	\item ¿Has podido orientarte correctamente (en un solo jugador) sin el sonar hasta poder completar el tutorial?
	\subitem- Sí, No, u otra respuesta de redacción libre.
	\item ¿Te has sentido confundido sobre lo que debías hacer durante el puzle?
	\subitem- \say{Sí, he necesitado la guía}, \say{No, lo he podido resolver sin la guía} u otra respuesta de redacción libre.
	\item Otras consideraciones y opiniones
	\subitem- Texto de respuesta larga.
	\begin{itemize}
		\item \textit{Preguntas sobre el puzle del piano.}
	\end{itemize}
	\item ¿Te ha parecido que esta mecánica de juego sonificada tiene potencial para ser utilizada en futuros juegos?
	\subitem- Respuesta del 1 al 10. 1 significa \say{Ha sido frustrante}, 10 significa \say{Ha sido divertida}.
	\item ¿Has podido ubicar correctamente los elementos del puzle del piano? (Melodía y el propio piano)
	\subitem- Sí, No, u otra respuesta de redacción libre.
	\item ¿El proceso de resolución ha sido intuitivo o has necesitado utilizar la guía para resolverlo?
	\subitem- \say{Ha sido intuitivo}, \say{He necesitado la guía} u otra respuesta de redacción libre.
	\item Otras consideraciones y opiniones
	\subitem- Texto de respuesta larga.
	\begin{itemize}
		\item \textit{Preguntas sobre el puzle de la pesca.}
	\end{itemize}
	\item ¿Te ha parecido que esta mecánica de juego sonificada tiene potencial para ser utilizada en futuros juegos?
	\subitem- Respuesta del 1 al 10. 1 significa \say{Ha sido frustrante}, 10 significa \say{Ha sido divertida}.
	\item ¿Has podido ubicar correctamente el puzle?
	\subitem- Sí, No, u otra respuesta de redacción libre.
	\item ¿La explicación dada por la voz es suficiente para entender cómo resolver el puzle?
	\subitem- \say{Sí, es suficiente}, \say{No, tuve que utilizar la guía} u otra respuesta de redacción libre.
	\item Otras consideraciones y opiniones
	\subitem- Texto de respuesta larga.
	\begin{itemize}
		\item \textit{Preguntas sobre el puzle de la pared y la llave.}
	\end{itemize}
	\item ¿Te ha parecido que esta mecánica de juego sonificada tiene potencial para ser utilizada en futuros juegos?
	\subitem- Respuesta del 1 al 10. 1 significa \say{Ha sido frustrante}, 10 significa \say{Ha sido divertida}.
	\item ¿Has podido ubicar correctamente el la llave  y la pared?
	\subitem- Sí, No, u otra respuesta de redacción libre.
	\item ¿El proceso de resolución ha sido intuitivo?
	\subitem- Sí, No, u otra respuesta de redacción libre.
	\item Otras consideraciones y opiniones
	\subitem- Texto de respuesta larga.
	\begin{itemize}
		\item \textit{Preguntas sobre el puzle de las salas.}
	\end{itemize}
	\item ¿Te ha parecido que esta mecánica de juego sonificada tiene potencial para ser utilizada en futuros juegos?
	\subitem- Respuesta del 1 al 10. 1 significa \say{Ha sido frustrante}, 10 significa \say{Ha sido divertida}.
	\item En el modo de un solo jugador, ¿es posible entender el funcionamiento del puzle con la información proporcionada?
	\subitem- Sí, No, u otra respuesta de redacción libre.
	\item En el modo de un solo jugador, ¿la interfaz de las teclas y su significado ha sido intuitiva?
	\subitem- Sí, No, u otra respuesta de redacción libre.
	\item Otras consideraciones de la versión de un solo jugador
	\subitem- Texto de respuesta larga.
	\item En la versión multijugador, ¿se entiende correctamente el papel de cada jugador en el puzle?
	\subitem- Sí, No, u otra respuesta de redacción libre.
	\item En la versión multijugador, ¿Ha sido intuitivo asociar los sonidos con las baldosas?
	\subitem- Sí, No, u otra respuesta de redacción libre.
	\item Otras consideraciones de la versión multijugador.
	\subitem- Texto de respuesta larga.
	\begin{itemize}
		\item \textit{Preguntas sobre el puzle del fin del juego.}
	\end{itemize}
	\item ¿Has podido ubicar correctamente el puzle y los elementos que lo resuelven?
	\subitem- Sí, No, u otra respuesta de redacción libre.
	\item Otras consideraciones y opiniones
	\subitem- Texto de respuesta larga.
	\item[$\blacksquare$] Las siguientes preguntas están relacionadas con cuestiones sobre la industria de los videojuegos y accesibilidad.
	\item ¿Consideras interesante la adaptación realizada de minijuegos tradicionalmente visuales a su versión sonificada?
	\subitem- Sí, No, u otra respuesta de redacción libre.
	\item ¿Qué orientación debería tomar la industria de los videojuegos en cuestiones de accesibilidad visual? (Crear juegos inherentemente accesibles, adaptar juegos existentes, crear herramientas externas...)
	\subitem- Texto de respuesta larga.
	\item ¿Consideras que la industria está avanzando en accesibilidad para personas ciegas?
	\subitem- Respuesta del 1 al 10. 1 significa \say{No avanzan nada}, 10 significa \say{Avanzan mucho y satisfactoriamente}.
	\item ¿Crees que existe suficiente información o formación sobre cómo diseñar juegos accesibles para ciegos?
	\subitem- \say{Sí, existe}, \say{No, no existe}, \say{No es suficiente, pero existe alguna} u otra respuesta de redacción libre.
	\item ¿Cuáles son las principales barreras para implementar accesibilidad?
	\subitem- Texto de respuesta larga.
	\item ¿Qué crees que necesitaría la industria para mejorar la accesibilidad para personas ciegas?
	\subitem- Texto de respuesta larga.
	\item ¿Piensas que la accesibilidad puede ser una ventaja competitiva en la industria de los videojuegos?
	\subitem- Sí, No, u otra respuesta de redacción libre.
\end{enumerate}